% IV. Програмна реализация. По предписания скелет за „програмна система“
% (REQUIREMENTS.md А.3): обща структура, йерархия на класове, коментар на изходния
% текст на програмата.
% Ориентировъчен обем: 12–14 стр.
%
% Числата тук (брой тестове, брой файлове по слой, имена на конфигурации) се ЧЕТАТ от
% хранилището към момента на предаване, не се преписват от този скелет — те растат.
\section{Програмна реализация}
\label{sec:implementation}

Разделът описва как проектните решения от Раздел~\ref{sec:design} са изразени в код:
организацията на изходния текст, конфигурациите за компилация, реализацията на всеки
слой и --- най-същественото за верността на резултата --- дисциплината на разработка,
която не позволява поведение да бъде измислено.

\subsection{Обща структура на изходния код}
\label{subsec:code_structure}

Изходният текст е разделен на публични заглавни файлове и реализация. Директориите
следват слоевете от §\ref{subsec:structure}: графични примитиви, ядро (договори и
обработчици), разположение, контроли, декларативен слой, услуги на устройството,
инициализация, анимации и платформени слоеве.

Едно отклонение от очакваното заслужава обяснение: обработчиците живеят заедно с
договорите в общата директория на ядрото, а не в собствена. Причината е, че двете се
изменят заедно --- добавянето на свойство към договор изисква ред в таблицата на
съответния обработчик --- и разделянето им би създало две места, които трябва да
се редактират синхронно, без нищо да го налага.

Съответствието директория\,--\,слой\,--\,обем е дадено в \tabref{tab:dirs}
(§\ref{subsec:structure}). Имената на пространствата съответстват на тези в оригинала
с преобразуване на стила, а всеки публичен заглавен файл носи в коментар пълното име на
типа, от който произхожда --- така изходният текст остава съпоставим при проверка, без да
се налага търсене по подобие.

\subsection{Организация на компилацията}
\label{subsec:build}

Компилацията е организирана като набор от независими конфигурации, всяка в собствено
дърво: слоят без екран, конфигурация със статичен анализ, три платформи със санитайзери за
поведение, за нишки и за неинициализирана памет, и по една платформа за всеки платформен
слой.

Разделянето поражда практическия проблем, че пълната проверка е всъщност N независими
компилации. Организацията му е част от реализацията, а не околност: конфигурацията без екран
се изпълнява първа, за да напълни кеша на компилатора, която платформата за статичен анализ
след това преизползва --- при еднакви флагове тя плаща само за самия анализ.
Компилациите са инкрементални по подразбиране; изчистването е изричен избор, а не
предпазна мярка.

Режимите на работа са два и смесването им е обичайната загуба на време.
\textbf{Бързият вътрешен цикъл} компилира една конфигурация инкрементално и изпълнява само
тестовете, чиито имена съвпадат с даден израз --- секунди. \textbf{Пълната проверка} минава
през всички конфигурации, включително санитайзерите и статичния анализ --- десетки минути.

Инструментите са нарочно два отделни, а не един с ключ, защото изборът влияе върху навика:
инструмент, който по подразбиране прави пълната проверка, се използва рядко, а инструмент,
който по подразбиране прави бързата, се използва след всяка промяна. Пълната се пази за
проверка преди предаване.

\subsection{Реализация на слоя за преобразуване}
\label{subsec:handler_impl}

\subsubsection{Йерархия на класовете}

\begin{figure}[H]
\centering
\begin{tikzpicture}[font=\footnotesize, node distance=0.6cm,
  n/.style={draw, rounded corners=2pt, align=center, inner sep=4pt, minimum height=0.62cm},
  ar/.style={-Latex, thick}]
  \node[n, fill=blue!6, text width=5.6cm] (i) {\code{i\_view\_handler}\\\scriptsize нешаблонен интерфейс --- полиморфната дръжка};
  \node[n, fill=green!10, text width=6.6cm, below=of i] (base)
     {\code{view\_handler<Derived, Virtual, Platform>}\\\scriptsize шаблонна база: жизнен цикъл, таблици, точки за разширение};
  \node[n, text width=3.0cm, below left=0.7cm and 0.6cm of base] (bh) {\code{button\_handler}\\\scriptsize таблици + create};
  \node[n, text width=3.0cm, below right=0.7cm and 0.6cm of base] (lh) {\code{label\_handler}};
  \node[n, fill=red!8, text width=2.5cm, below=0.75cm of bh, xshift=-0.6cm] (p1) {\code{ios/}\\\scriptsize \code{.mm}};
  \node[n, fill=red!8, text width=2.5cm, below=0.75cm of bh, xshift=2.4cm] (p2) {\code{android/}};
  \node[n, fill=gray!10, text width=2.5cm, below=0.75cm of bh, xshift=5.4cm] (p3) {\code{headless/}};
  \draw[ar] (base) -- node[right,font=\scriptsize] {virtual} (i);
  \draw[ar] (bh) -- (base); \draw[ar] (lh) -- (base);
  \draw[ar] (p1) -- (bh); \draw[ar] (p2) -- (bh); \draw[ar] (p3) -- (bh);
\end{tikzpicture}
\caption{Йерархия на обработчиците. Наследяването на интерфейса е \emph{виртуално},
за да може обработчик на контейнер да удовлетвори едновременно него и интерфейса за
контейнери без два подобекта. Платформените реализации са отделни единици за превод, а не
класове --- добавянето на платформа е добавяне на файлове.}
\label{fig:handler_hierarchy}
\end{figure}

Разделянето на частичен клас с платформени файлове, което оригиналът използва, е
изразено чрез отделни единици за превод: обща реализация плюс по една реализация на
платформен слой, избирана при конфигуриране. Така добавянето на платформа е
\emph{добавяне} на файлове, а не изменение на съществуващи --- което е и проверката
дали границата между преносимото и платформеното е на правилното място.

\subsubsection{Коментар на изходния текст: един представителен срез}

Срезът за бутон е представителен --- всяка друга контрола е със същата форма ---
затова се разглежда подробно веднъж:

\begin{lstlisting}[caption={Таблицата за свойства на обработчика за бутон},label={lst:button_handler}]
// Текст и изглед на текста, ключирани по i_text_button.
property_mapper<i_text_button, button_handler>& button_handler::text_mapper()
{
    static property_mapper<i_text_button, button_handler> table{
        {"text",              &button_handler::map_text},
        {"text_color",        &button_handler::map_text_color},
        {"font",              &button_handler::map_font},
        {"character_spacing", &button_handler::map_character_spacing},
    };
    return table;
}
\end{lstlisting}

Четири неща в тези редове носят проектните решения от Раздел~\ref{sec:design}. Таблицата е
\code{static} --- изгражда се веднъж за целия процес, не за всеки бутон. Ключовете са
низове, защото са същите имена, които системата от свойства известява при промяна.
Стойностите са указатели към функции-членове, а не към виртуални методи --- разрешаването
е по време на компилация. И типът е шаблон по \emph{договора} (\code{i\_text\_button}), а
не по контролата, така че една и съща таблица обслужва всеки тип, който удовлетворява
договора.

Платформената страна допълва картината:

\begin{lstlisting}[caption={Създаване на native контролата (iOS)},label={lst:button_ios}]
std::unique_ptr<button_platform> button_handler::create_platform_view()
{
    auto platform = std::make_unique<button_platform>();
    UIButton* const button = [UIButton buttonWithType:UIButtonTypeSystem];
    set_control_properties_from_proxy(button);
    platform->native = (__bridge_retained void*)button; // слотът поема една препратка
    return platform;
}
\end{lstlisting}

Тук се вижда границата между двата модела на живот: \code{\_\_bridge\_retained} прехвърля
една препратка от управлението на средата към C++ слота, а разрушителят я връща обратно.
Смесеният език позволява това да се случи в \emph{една} функция --- без междинно
представяне и без слой за маршалиране.

% \begin{lstlisting}[caption={Реализация на обработчик}, label={lst:button_handler}]
% \end{lstlisting}

\subsection{Реализация на алгоритмите за разположение}
\label{subsec:layout_impl}

Разпределянето на пропорционалните измерения е с порядък по-обемно от всеки друг
мениджър на разположение --- само по себе си довод за решението алгоритмите да се пренасят
водени от тестове, а не да се пишат по описание: алгоритъм с толкова разклонения има
твърде много гранични случаи, за да бъдат отгатнати.

\begin{lstlisting}[caption={Сумиране на теглата на пропорционалните измерения},label={lst:grid_stars}]
double star_count = 0.0;
for (const definition& current : definitions)
{
    if (current.is_star())
    {
        star_count += current.length().value();
    }
}
return star_count;
\end{lstlisting}

Функцията изглежда тривиална и точно затова е показателна: сумата се смята върху
\emph{всички} определения, а не само върху незаетите, защото клетка, разпростряна върху
няколко колони, вече е внесла своя дял. Разделянето на остатъка после става на
\code{остатък / star\_count} на единица тегло. Грешка тук --- например изключване на
колона, чиято ширина е вече разрешена --- не се проявява при една пропорционална колона и
се проявява при две с различни тегла, което е и причината този алгоритъм да се пренася
воден от тестове.

\subsection{Реализация на системата от свойства}
\label{subsec:binding_impl}

Стойностите се съхраняват не като една стойност, а като \emph{подредена по предимство}
редица: всеки източник (подразбиращ се, стил, свързване, ръчно присвояване, тригер,
визуално състояние, обработчик) заема свое ниво, а действащата стойност е тази с
най-високото заето ниво. Изчистването на един източник не изтрива стойността, а разкрива
следващата отдолу --- което е точното поведение, което оригиналът има.

Стойността, зададена от обработчика, се съхранява на най-високото ниво, но се
измества от всяко друго присвояване. Причината е, че тя не е избор на приложението, а
отражение на това, което native контролата е направила сама --- и всяко изрично
присвояване трябва да я надделее.

Разрешаването на пътя при свързване се извършва \emph{без рефлексия}: всяко свързваемо
свойство се вписва в регистър при своята регистрация, а изразът за свързване се
разрешава срещу този регистър.

\subsection{Реализация на платформените слоеве}
\label{subsec:platform_impl}

Петте платформени слоя се различават не само по инструментариум, но и по езиковата
граница, през която минават.

\subsubsection{Слоевете върху Apple (macOS и iOS)}

Слоевете върху Apple се пишат на смесен език (\code{.mm}), при което една единица за превод
вижда едновременно C++ и native обектния модел --- native обектът се държи пряко в C++ клас,
вместо да се пренася през граница (листинг~\ref{lst:app_bh_ios}). Управлението на живота се
съгласува изрично при преминаване на слота и обратно при разрушаване.

\paragraph{Защо това не превръща целия проект в смесен език.} Естественото притеснение е, че
щом част от файловете се компилират като Objective-C++, това ще плъзне нагоре: заглавните
файлове ще започнат да носят native типове и всеки консуматор ще трябва да се компилира по
същия начин. Не се случва, и причината е една --- \emph{native типът никъде не се появява в
заглавен файл}. Платформената структура държи native обекта зад непрозрачен указател:

\begin{lstlisting}[caption={Native обектът зад непрозрачен слот},label={lst:opaque},
  basicstyle=\scriptsize\ttfamily]
struct button_platform : view_platform_base
{
    void* native = nullptr;   // UIButton* / NSButton* — типът е известен само на .mm
};
\end{lstlisting}

Заглавният файл вижда \code{void*} и нищо повече. Само реализацията --- която \emph{вече} е
\code{.mm} --- знае какъв е действителният тип и извършва преобразуването. Следствията са
три: кросплатформените заглавни файлове се компилират от обикновен C++ компилатор; тестовете
без екран не се нуждаят от Apple инструментариум; и смяна на native типа не предизвиква
повторна компилация извън собствената му единица за превод.

\paragraph{Как изграждането отделя платформената част.} Разделянето не е чрез условна
компилация вътре във файловете, а чрез \emph{подбор на файлове}. Изграждането има една
променлива --- избраният платформен слой --- и според нея включва точно една директория:

\begin{center}
\begin{tabular}{@{}L{2.8cm}L{4.0cm}L{6.8cm}@{}}
\hline
\textbf{Стойност} & \textbf{Директория} & \textbf{Език на единиците} \\
\hline
\code{headless} & \code{platform/headless/} & обикновен C++ \\
\code{apple}    & \code{platform/apple/}    & Objective-C++ (AppKit) \\
\code{ios}      & \code{platform/ios/}      & Objective-C++ (UIKit) \\
\code{maccatalyst} & \code{platform/ios/}   & същите файлове, друга целева платформа \\
\code{windows}  & \code{platform/windows/}  & C++ с проекция на компонентния модел \\
\code{android}  & \code{platform/android/}  & C++ през интерфейса за извикване \\
\hline
\end{tabular}
\end{center}

Два реда заслужават внимание. \code{maccatalyst} не е отделен слой --- изграждането
пренасочва към \emph{същите} файлове на UIKit и само сменя целевата платформа; това е и
причината настолната Apple платформа да получи два различни слоя без двойно писане
(§\ref{subsec:backends}). А \code{headless} е обикновен C++ и не изисква нито един платформен
инструментариум --- затова проверката на преносимата логика се изпълнява навсякъде, включително
там, където Apple компилатор няма.

Онова, което наистина пази разделението, е правилото от §\ref{subsec:structure}: зависимостите
сочат само надолу. Ако кросплатформен заглавен файл можеше да достигне до платформен тип, нито
подборът на файлове, нито непрозрачният указател биха помогнали --- разделението щеше да бъде
обходено още при включването.

\subsubsection{Слоят върху Android}
Слоят върху Android минава през интерфейса за извикване към управляваната платформа.
Той е най-скъпият по К5 от §\ref{subsec:criteria} по три причини: всяко обръщение изисква
търсене на метод по подпис; препратките към обекти имат два вида живот --- локален (валиден
до връщане от функцията) и глобален (изрично освобождаван), и смесването им е източник на
дефекти, които се проявяват късно; и всяка нишка, която обръща към платформата, трябва да
бъде изрично присъединена и отделена.

Жизненият цикъл на дейността се съгласува с този на приложението чрез изрично препредаване
на събитията --- платформата може да разруши и пресъздаде дейността, без приложението да е
приключило.

\subsubsection{Слоят върху Windows}
Слоят върху Windows минава през проекция на компонентния модел в заглавни файлове.
Основното ограничение, което инструментариумът налага, е върху \emph{свободното
разположение}: контейнерът, който позволява произволни координати, се държи различно от
този, който подрежда сам, и системата трябва да избере правилния според това какво
изисква описанието.

Отбелязват се и три езикови капана, установени в хода на работата, защото всеки от тях
струва часове и нито един не е документиран като такъв: имена от системните заглавни
файлове, които съвпадат с имена от системата (разрешава се с преименуване или с
премахване на макрос); съкратено име на пространство, което влиза в конфликт с име на
тип; и невъзможността да се преопредели заделянето на памет за проектираните типове.

Обхватът на реализираните контроли на тази платформа е \emph{частичен по замисъл} към
момента на предаване: реализирани са прозорецът, страницата, контейнерите за разположение,
надписът и бутонът, а останалите контроли все още използват огледалото без екран и не
изобразяват нищо. Това се отбелязва изрично тук, защото иначе би се подразбрало от
таблиците --- интерфейс, изграден само от нереализирани контроли, е празен и в двете
колони, което е незавършена реализация, а не съвпадение.

\subsubsection{Слоят без екран}
Слоят без екран реализира същите обработчици, но „native контролата“ е обикновена
структура, която само записва получените стойности. Това го прави пълноценен обект на
автоматизирана проверка: всичко, което не е рисуване --- разположение, предимство на
стойностите, известяване, жизнен цикъл --- е проверимо върху него, без устройство и без
сравняване на изображения.

Той е и \emph{определението} за това какво трябва да реализира един платформен слой:
новият слой се сверява срещу него, а не срещу някой от вече съществуващите, които носят
собствените си особености.

\subsection{Реализация на декларативния слой}
\label{subsec:xaml_impl}

\subsubsection{Входът: едно описание, два потребителя}

Описанието е обикновен XML документ. Съществено е, че \emph{същите байтове} се четат и от
еталонното приложение, и от системата --- това е условието, което прави сравнението в
Раздел~\ref{sec:method} валидно.

\begin{lstlisting}[caption={Откъс от описание на интерфейс},label={lst:xaml_src},
  basicstyle=\scriptsize\ttfamily]
<ContentPage xmlns="http://schemas.microsoft.com/dotnet/2021/maui"
             xmlns:x="http://schemas.microsoft.com/winfx/2009/xaml"
             x:Class="MauiReference.Pages.BoxViewPage" Title="BoxView">
    <ScrollView>
        <VerticalStackLayout Spacing="6" Padding="12">
            <Label Text="Default" />
            <BoxView BackgroundColor="CornflowerBlue"
                     WidthRequest="160" HeightRequest="160" />
        </VerticalStackLayout>
    </ScrollView>
</ContentPage>
\end{lstlisting}

\subsubsection{Конвейерът}

Описанието влиза в единицата за превод чрез \code{\#embed} --- директива на C++26, която
вгражда съдържанието на файл като масив от байтове по време на компилация. Оттам нататък
всичко е обикновена \code{constexpr} обработка:

\begin{lstlisting}[caption={Вграждане и изграждане на страница по време на компилация},
  label={lst:xaml_build}, basicstyle=\scriptsize\ttfamily]
namespace
{
    // Суровият markup, вграден от компилатора (пътят е спрямо този файл).
    constexpr unsigned char data_binding_xaml_bytes[] = {
        #embed "data_binding.xaml"
    };
    constexpr maui::fixed_string data_binding_xaml{data_binding_xaml_bytes};
}

std::unique_ptr<maui::controls::content_page> data_binding_page(
        std::unique_ptr<ViewModels::greeting_view_model> vm)
{
    auto page = maui::build_page<ViewModels::greeting_view_model, data_binding_xaml>();
    page->bind_to(std::move(vm)); // BindingContext -> {Binding} изразите оживяват
    return page;
}
\end{lstlisting}

Ключът е сигнатурата \code{build\_page<ViewModel, fixed\_string>()}: описанието е
\emph{шаблонен параметър}, не аргумент. Това го прави достъпно за компилатора, който
разбира документа и изгражда съответните извиквания на конструктори --- без разбор при
стартиране и без нито един ред кодогенерация извън езика: няма отделна стъпка в изграждането
и няма породени файлове в дървото.

Трите стъпки от §\ref{subsec:xaml_design} се разпределят така:

\begin{center}
\begin{tabular}{@{}L{3.4cm}L{3.0cm}L{7.0cm}@{}}
\hline
\textbf{Стъпка} & \textbf{Кога} & \textbf{Какво изисква} \\
\hline
разбор & при компилация & \code{constexpr} обработка на вградените байтове \\
изграждане на графа & при компилация & регистър на типовете (замества рефлексията) \\
свързване на данни & при изпълнение & зависи от данните, не от описанието \\
\hline
\end{tabular}
\end{center}

Печалбата е в момента, в който грешката се открива: непознат елемент или сгрешено име на
свойство спира \emph{компилацията}, вместо да предизвика отказ при стартиране.

\subsubsection{Границата на средството}

\code{\#embed} не е налична във всички компилатори, които проектът поддържа. Там, където
липсва, същият вход минава през отделна стъпка в изграждането, която го превръща в масив от
байтове --- еднакъв резултат, една допълнителна зависимост за тази платформа. Цената на целия
декларативен път е измерена в §\ref{subsec:cost} и се оказва \emph{код}, а не данни.

\subsection{Обработка на събития от интерфейса}
\label{subsec:events_impl}

Дотук стойностите текат в една посока --- от приложението към native контролата. Събитията
текат в обратната и изискват собствен механизъм, защото native страната е тази, която ги
поражда.

\subsubsection{Механизмът от страната на C++}

Няма делегати и няма събиране на боклука, затова събитието е самостоятелен тип:
\code{event<Args...>} с \code{connect} (връща \term{жетон}), \code{disconnect} и
\code{raise}. Жетонът е съществен --- той позволява абонаментът да бъде прекратен
\emph{детерминистично}, а не при следващо събиране. Има и обвивка, която прекратява
абонамента в разрушителя си, така че животът на връзката съвпада с обхвата, в който е
създадена.

\subsubsection{Пътят от native контролата до приложението}

Между двата обектни модела стои малък посредник. На Apple платформите това е обект от native
типовата система, който няма друга задача освен да препрати:

\begin{lstlisting}[caption={Посредник, който препраща native действието към виртуалния изглед},
  label={lst:event_proxy}, basicstyle=\scriptsize\ttfamily]
// Obj-C trampoline: препраща target-action на UIButton към виртуалния изглед.
@interface MauiButtonEventProxy : NSObject
@property(nonatomic) maui::core::button_handler* handler;
- (void)onTouchUpInside:(id)sender;
@end
\end{lstlisting}

Веригата е: native действие $\rightarrow$ посредник $\rightarrow$ обработчик $\rightarrow$
\code{send\_clicked()} на виртуалния изглед $\rightarrow$ контролата издига собственото си
събитие $\rightarrow$ обработчиците на приложението.

Два детайла не са произволни. Първо, \emph{редът} на събитията се пренася точно: при
отпускане на пръста върху бутона първо се издига „отпуснат“ и едва след това „натиснат“ ---
редът на еталона, който приложение може да наблюдава. Второ, посредникът държи
\emph{непритежаваща} препратка към обработчика; притежаваща би затворила цикъл (обработчикът
притежава native контролата, която държи посредника) --- точно класът цикли, който моделът
на собственост от §\ref{subsec:ownership} прави невъзможен по-нагоре в стека.

\subsubsection{Именувани събития за декларативния слой}

Описанието може да свърже действие по \emph{име} (например атрибут \code{Clicked}). Тъй като
рефлексия няма, всяка контрола вписва събитията си в регистър при изграждането си ---
съответствие „име $\rightarrow$ функция, която абонира“. Декларативният слой търси там,
намира функцията и я извиква. Същата размяна както при свойствата: явен запис вместо
откриване по време на изпълнение.

\subsection{Дисциплина на разработка и осигуряване на верността}
\label{subsec:process}

Този параграф описва процеса, а не кода --- и е поставен в раздела за реализация
съзнателно, защото при работа с външен еталон именно процесът е това, което отличава
пренасяне на поведение от повторното му измисляне.

\subsubsection{Цикълът „по една компонента“}

За всяка компонента редът е фиксиран: прочитане на договора и на \emph{тестовете} на
оригинала; пренасяне на тестовете \emph{преди} реализацията, така че те да се
компилират и да се провалят; реализация до преминаването им; и чак тогава преминаване
към следващата. Правилото, което този ред налага, е формулирано отрицателно и е
по-важно от самия ред: \emph{когато поведението не е известно, то се прочита от
оригинала --- не се предполага}. Правдоподобно предположение, което се разминава с
действителното поведение, е дефект, а не приблизително решение.

Трите нива на истината и въпросите, на които отговарят, са изложени в
§\ref{subsec:port_strategy}: описанието на публичната повърхност дава \emph{формата},
тестовете дават \emph{поведението}, изходният текст дава \emph{реализацията}.

Разминаването между тях е реално и е било разрешавано. Пример: изходният текст на еталона
поставя долна граница на размера на една контрола; разпространената версия изобразява
по-малък размер без никаква граница. Формата (публичната повърхност) не казва нищо по
въпроса, тестовете също. Приложеното правило --- предимство има изобразеното --- е записано
предварително (§\ref{subsec:rulings}) и промяната е отбелязана в кода като документирано
отклонение с цитиране на двата източника. Без предварително записано правило това решение
би било произволно.

\subsubsection{Статичен анализ и санитайзери като част от цикъла}

Статичният анализ се изпълнява като отделна конфигурация с изискване за \emph{нула}
находки, и потискането на находка вместо отстраняването ѝ не е допустимо. Санитайзерите ---
за неопределено поведение, за състезания между нишки и за неинициализирана памет --- са три
отделни конфигурации, защото не могат да се съчетаят в една.

Всичко това е част от цикъла, а не заключителна проверка, по проста сметка: находка,
открита седмици след внасянето си, изисква възстановяване на контекста, който вече не е в
главата на никого, докато същата находка в цикъла, в който е внесена, се отстранява
веднага. Разликата е в порядъци, не в проценти.

\subsubsection{Обхват на автоматизираната проверка}

\begin{table}[H]
\caption{Обем на автоматизираната проверка по слой}
\label{tab:tests}
\centering\small
\begin{tabular}{@{}L{4.2cm}rL{6.4cm}@{}}
\hline
\textbf{Слой} & \textbf{Файлове} & \textbf{Какво покриват} \\
\hline
графични примитиви & 13  & стойностни типове, геометрия \\
ядро (договори, таблици) & 17 & свойства, предимство, жизнен цикъл \\
разположение & 9 & измерване и подреждане на всеки мениджър \\
контроли & 190 & поведението на всяка контрола \\
декларативен слой & 9 & разбор, изграждане, свързване \\
услуги на устройството & 34 & — \\
инициализация & 4 & регистрация, изграждане на приложението \\
анимации & 5 & — \\
платформени слоеве & 19 & връзките към native страната \\
\hline
\end{tabular}
\end{table}

По записа от дневника на разработката към последната пълна проверка конфигурацията без
екран изпълнява \textbf{3765} случая, всички преминаващи. Числото расте с всяка добавена
компонента и се чете наново непосредствено преди предаване.

Съществено е \emph{какво този обем не доказва}. Той измерва преносимата част: разположение,
свойства, известяване, жизнен цикъл. Изобразяването не се проверява от него --- то се
проверява от методиката в Раздел~\ref{sec:method}. Двете мерки са допълващи се и нито една
не замества другата: система с изцяло преминаващи тестове може да изобразява погрешно, а
система с точно съвпадение на изображенията може да има грешна семантика, която тези
интерфейси просто не докосват.
